\newpage
\section{Spanning Sets and Linear Independence}
Consider a vector space $V = \mathbb{R}^3$ and three vectors $\hat{\imath}, \hat{\jmath}$, and $\hat{k}$ where $\hat{\imath} = (1, 0, 0), \hat{\jmath} = (0, 1, 0), \hat{k} = (0, 0, 1) \in \mathbb{R}^3$. Then every vector $(a, b, c) \in \mathbb{R}^3$ can be expressed as 

$$(a, b, c) = (a, 0, 0) + (0, b, 0) + (0, 0, c)$$
$$=a\hat{\imath} + b\hat{\jmath} + c\hat{k}$$

Every vector in $\mathbb{R}^3$ can be expressed in terms of $\hat{\imath}, \hat{\jmath}$, and $\hat{k}$.

\begin{definition}
Let $V$ be a vector space and $v_1, v_2, \dots, v_k \in V$. A \textbf{linear combination} of $v_1, v_2, \dots, v_k$ is any element of the form
$$a_1v_1 + a_2v_2 + \cdots + a_kv_k = \sum_{i = 1}^k a_iv_i$$
where $a_i \in \mathbb{R}$ for all $i$.
\end{definition}

\begin{example} The following are examples of linear combinations:

\begin{enumerate}
\item $0v_1 + 0v_2 + \cdots + 0v_k = \vec{0}$ is a linear combination of $v_1, v_2, \dots, v_k$.

\item Every vector in $\mathbb{R}^3$ is a linear combination of $\hat{\imath}, \hat{\jmath}, \hat{k}$.

\item In the vector space $P_4$, $4x^4 - 2x^3 + 6x - 9 = 4x^4 + (-2)x^3 + 0x^2 + 6x + (-9)(1)$ is a linear combination of $1, x, x^2, x^3, x^4$.

\item Let $A = \bmat{1 & -3 & 0 \\ 0 & 2 & 2}$ and $x = \bmat{a \\ b \\ c}$. Then
\begin{align*}
Ax = \bmat{1 & -3 & 0 \\ 0 & 2 & 2}\bmat{a \\ b \\ c} = \bmat{a - 3b \\ 2b + 2c} \\
=\bmat{a \\ 0} + \bmat{-3b \\ 2b} + \bmat{0 \\ c} \\
= a\bmat{1 \\ 0} + b\bmat{-3 \\ 2} + c\bmat{0 \\ 2}
\end{align*}

$Ax$ is a linear combination of the columns of $A$. In general if $A \in M_{m \times n}(\mathbb{R})$ and $x \in M_{n \times 1}(\mathbb{R})$, then $Ax$ is a linear combination of the columns of $A$.

\item Let $v = (-3, -\frac{1}{2}, \frac{5}{2})$. Is $v$ a linear combination of $(1, 0, 0), (2, 1, 0)$, and $(0, 1, 1)$?

\textit{Solution.} We want to know if there exists scalars $a, b, c$ such that
$$\left(-3, -\frac{1}{2}, \frac{5}{2}\right) = a\left(1, 0, 0\right) + b(2, 1, 0) + c(0, 1, 1)$$
$$=(a + 2b, b + c, c)$$

Comparing coordinates,

\item In $V = \mathscr{F}(\mathbb{R})$, $\cosh \in V$ is a linear combination of $e^x$ and $e^{-x}$
$$\cosh x = \frac{e^x + e^{-x}}{2} = \frac{1}{2}e^x + \frac{1}{2}e^{-x}$$
\end{enumerate}


\end{example}

\begin{definition}
Let $S = \{v_1, v_2, \dots, v_k\}$ be a nonempty subset of a vector space $V$
\begin{itemize}
    \item The span of $S$, denoted by $\operatorname{span} S$, is the set of all linear combinations of elements of $S$, that is,
    $$\operatorname{span} S = \{a_1v_1 + a_2v_2 + \cdots + a_kv_k \mid a_i \in \mathbb{R}\,\, \forall i\}$$
    \item If $\operatorname{span} S = V$, we say that $S$ is a \underline{spanning set} of $V$ or $V$ is spanned by $S$.\footnote{Every vector in $V$ is a linear combination of elements in $S$.}
\end{itemize}
\end{definition}

\begin{example}
Let $V = \mathbb{R}^2$ and $S = \{(1, 0, 0), (0, 1, 0), (0, 0, 1)\}$, then,
$$\operatorname{span}S = \{a(1,0, 0) + b(0, 1, 0) + c(0, 0, 1) \mid a, b, c \in \mathbb{R}\}$$
$$=\{(a, b, c) \mid a, b, c \in \mathbb{R}\}$$
Therefore, $S$ spans $\mathbb{R}^3$.
\end{example}

\begin{example}
Let $V = P_n$ and $S = \{1, x, x^2, \dots, x^n\}$, then,
$$\operatorname{span} S = \{a_0 + a_1x + a_2x^2 + \cdots + a_nx^n \mid a_i \in \RR\,\, \forall i\} = P_n$$
Therefore, $S$ spans $P_n$.
\end{example}

\begin{example}
Let $V = \CC = \{a + bi \mid a, b \in \RR\,\,\text{and}\,\,i^2 = -1\}$ and $S = \{1, i\}$, then,
$$\operatorname{span}S = \{a + bi \mid a, b \in \RR\,\,\text{and}\,\,i^2 = -1\} = \CC$$
Therefore, $S$ spans $\CC$.
\end{example}

\begin{example}
Let $V = M_2(\RR)$ and $S = \left\{\bmat{1 & 0 \\ 0 & 0}, \bmat{ 0 & 1 \\ 0 & 0}, \bmat{0 & 0 \\ 1 & 0}, \bmat{0 & 0 \\ 0 & 1}\right\}$, then,
$$\operatorname{span}S = \left\{\bmat{a & b \\ c & d} \suchthat a,b, c, d \in \RR \right\} = M_2(\RR)$$
Therefore $S$ spans $M_2(\RR)$.
\end{example}

\begin{example}
Let $S = \{(2, 3), (1, 2)\}$. Does this span $\RR^2$?
\end{example}

\textit{Solution.} Let $(a, b) \in \RR^2$. We want to know if there exists $x, y \in \RR$ such that,
$$(a, b) = x(2, 3) + y(1, 2)$$
$$=(2x, 3x) + (y, 2y)$$
$$=(2x + y, 3x + 2y)$$
\begin{align*}
\begin{cases}
    2x + y = a \\
    3x + 2y = b 
\end{cases}
\quad
\longleftrightarrow
\quad
\bmat{2 & 1 \\ 3 & 2}\bmat{x \\ y} = \bmat{a \\ b}
\end{align*}
The coefficient matrix $A$ of the system is nonsingular, since $\det A \neq 0$, therefore the system has a solution.
$$\bmat{x \\ y} = \bmat{2 & 1 \\ 3 & 2}\inv \bmat{a \\ b}$$
$$= \bmat{2 & -1 \\ 3 & -2}\bmat{a \\ b} = \bmat{2a - b \\ 3a - 2b}$$
Therefore, $S$ spans $\RR^2$.

\newpage

\begin{example}
Let $S = \{x^2 + 2x + 1, x^2 + 2\} \subseteq P_2$. Does $S$ span $P_2$?
\end{example}

\textit{Solution.} Let $ax^2 + bx + c \in P_2$. We want to know if there exists $D, E \in \RR$ such that
$$ax^2 + bx + c = D(x^2 + 2x + 1) + E(x^2 + 2)$$
$$=x^2(D + E) + x(2D) + (D + 2E)$$
\begin{align*}
\begin{cases}
    D + E = a \\ 
    2D = b \\
    D + 2E = c
\end{cases}
\quad
\longleftrightarrow
\quad 
\begin{amatrix}{2}
    1 & 0 & a \\
    2 & 0 & b \\
    1 & 2 & c
\end{amatrix}
\xrightarrow{EROs}
\begin{amatrix}{2}
    1 & 1 & a \\
    0 & 1 & -a+ c \\
    0 & 0 & -4a + b + 2c \\
\end{amatrix}
\end{align*}
Since there is a restriction placed on the last row, such that $-4a + b + 2c$ must be equal to 0, therefore $S$ does not span $P_2$.

\vspace{0.2cm}

\begin{theorem}
If $S = \{v_1, v_2, \dots, v_k\}$ is a nonempty subset of $V$, then $\operatorname{span}S \leq V$.
\end{theorem}

\begin{remark}Shown below are some remarks regarding the span of sets:
\begin{itemize}
    \item If $S \subseteq \operatorname{span} S$ and $\operatorname{span}S$ is the smallest subspace of $V$ containing $S$.
    \item $\operatorname{span}(\emptyset) = \{\vec{0}\}$
\end{itemize}
\end{remark}

\begin{definition}
Let $S = \{v_1, v_2, \dots, v_n\}$ be a nonempty subset of a vector space $V$. Then $S$ is said to be \textbf{linearly independent} if whenever
$$a_1v_1 + a_2v_2 + \cdots + a_nv_n = \vec{0}$$
then $a_1 = a_2 = \cdots = a_n = 0$. Otherwise, $S$ is said to be \textbf{linearly dependent}.
\end{definition}

\begin{remark}
The following are some remarks on linear dependence:
\begin{enumerate}
    \item $S$ is linearly dependent iff there exists $a_1, a_2, \dots, a_n \in \RR$ and not all zero, such that,
    $$\sum_{i=1}^na_iv_i = \vec{0}$$
    \item $\emptyset$ is taken to be linearly independent.
\end{enumerate}
\end{remark}

\begin{example}\label{example:6.5}
Determine if the set $S = \{(1, 0, 0), (0, 1, 0), (0, 0, 1)\} \subseteq \mathbb{R}^3$ is linearly independent.

\textit{Solution.} Form the equation,
$$a(1, 0, 0) + b(0, 1, 0) + c(0, 0, 1) = (0, 0, 0)$$
$$(a, b, c) = (0, 0, 0)$$
Because $a = b = c = 0$, therefore $S$ is linearly independent.

In general, $S = \{e_1, e_2, \dots, e_n\}$ is a linear independent subset of $\RR^n$.
\end{example}

\begin{example}\label{example:6.6}
Determine whether the set $S = \left\{\bmat{1 & 0 \\ 0 & 0}, \bmat{0 & 1 \\ 0 & 0}, \bmat{0 & 0 \\ 1 & 0}, \bmat{0 & 0 \\ 0 & 1}\right\} \subseteq M_2(\RR)$ is linearly independent.

\textit{Solution.} Form the equation,
$$a\bmat{1 & 0 \\ 0 & 0} + b\bmat{0 & 1 \\ 0 & 0} + c\bmat{0 & 0 \\ 1 & 0} + d\bmat{0 & 0 \\ 0 & 1} = \bmat{0 & 0 \\ 0 & 0}$$
$$=\bmat{a & 0 \\ 0 & 0} + \bmat{0 & b \\ 0 & 0} + \bmat{0 & 0 \\ c & 0} + \bmat{0 & 0 \\ 0 & d} = \bmat{0 & 0 \\ 0 & 0}$$
$$=\bmat{a & b \\ c & d} = \bmat{0 & 0 \\ 0 & 0}$$
where $a, b, c, d \in \RR$. Since $a = b = c = d = 0$, therefore, $S$ is linearly independent.
\end{example}

\lline

In general, $S = \{E_{k\ell} \mid k = 1, 2, \dots, m, \quad \ell = 1, 2, \dots, n\}$ is a linear independent subset of $M_{m \times n}(\RR)$. As an exercise, determine if the set $S$ is a spanning set in \hyperref[example:6.5]{Example 6.5} and \hyperref[example:6.6]{Example 6.6}.

\begin{example}
Determine if the set $S = \{e^x, e^{-x}\} \subseteq \mathscr{F}(\RR)$.

\textit{Solution.} Form the equation,
$$ae^x + be^{-x} = 0$$
where $a, b \in \RR$. By differentiating both sides we obtain,
$$ae^x -be^{-x} = 0$$
Adding the first and second equation together,
$$ae^x + be^{-x} + ae^x - be^{-x} = 0$$
$$2ae^x = 0 \Longrightarrow a = 0$$
Since $a = b = 0$, therefore, $S$ is linearly independent.
\end{example}

\begin{remark}
Let $V$ be a vector space and $S, T \leq V$ be finite:
\begin{enumerate}
    \item If $\vec{0} \in S$, then $S$ is linearly dependent.
    \item If $\vec{0} + v \in V$, then $\{v\}$ is linearly independent.
    \item Suppose $S \subseteq T$
    \begin{enumerate}
        \item $\operatorname{span}S \subseteq \operatorname{span}T$
        \item If $T$ is linearly independent then $S$ is linearly independent
        \item If $S$ is linearly dependent then $T$ is linearly dependent
    \end{enumerate}
    \item $S$ is linearly dependent iff some vector in $S$ is a linear combination of the other vectors in $S$. In particular, $S = \{u, v\}\quad(u, v \ne 0)$ is linearly dependent iff $v = cu$ for some $0 \ne c \in \RR$. In this case, $u$ and $v$ are said to be parallel.
\end{enumerate}
\end{remark}